\documentclass[10pt,a4paper]{article}
\usepackage[T1]{fontenc}
\usepackage{lmodern}
\usepackage[margin=23mm]{geometry}
\usepackage{amsmath,amssymb,amsthm}
\usepackage{xcolor,microtype}
\definecolor{assumption}{RGB}{155,65,0}
\newcommand{\extra}[1]{\textcolor{assumption}{#1}}
\newcommand{\op}[1]{\mathcal{#1}}
\newtheorem{lemma}{Lemma}
\renewcommand{\thelemma}{4.\arabic{lemma}}
\setcounter{lemma}{13}
\setlength{\parindent}{0pt}
\setlength{\parskip}{5pt}
\allowdisplaybreaks
% Adapted from Scherer and Weiland, Linear Matrix Inequalities in Control,
% November 2016 compilation, Lemmas 4.14--4.15, pp. 111--113.
\begin{document}

\begin{lemma}[Projection Lemma]
Let $\op P=\op P^*$ be a bounded PI operator partitioned into three
rows/columns and consider the operator inequality
\begin{equation}\tag{4.4.5}\label{eq:445}
 \begin{pmatrix}
 \op P_{11}&\op P_{12}+\op X^*&\op P_{13}\\
 \op P_{21}+\op X&\op P_{22}&\op P_{23}\\
 \op P_{31}&\op P_{32}&\op P_{33}
 \end{pmatrix}\prec0
\end{equation}
in the unstructured PI operator $\op X$. Here and below, a self-adjoint
operator $\op R$ is \emph{coercive} if $\op R\succeq\varepsilon I$
for some $\varepsilon>0$, and $\op L\prec0$ means that $-\op L$
is coercive. We use the PI algebra, which is closed under inversion
of coercive operators.
There exists a solution $\op X$ of \eqref{eq:445} if and only if
\begin{equation}\tag{4.4.6}\label{eq:446}
 \begin{pmatrix}\op P_{11}&\op P_{13}\\
                 \op P_{31}&\op P_{33}\end{pmatrix}\prec0
 \quad\text{and}\quad
 \begin{pmatrix}\op P_{22}&\op P_{23}\\
                 \op P_{32}&\op P_{33}\end{pmatrix}\prec0.
\end{equation}
If \eqref{eq:446} hold, one particular solution is given by
\begin{equation}\tag{4.4.7}\label{eq:447}
 \op X=\op P_{32}^*\op P_{33}^{-1}\op P_{31}-\op P_{21}.
\end{equation}
\end{lemma}

\begin{proof}
(Only if) If \eqref{eq:445} has a solution then \eqref{eq:446} follow
by canceling the first or second block row/column.

(If) Suppose that \eqref{eq:446} holds. Then $-\op P_{33}$ is coercive,
so $\op P_{33}^{-1}$ is a bounded PI operator. The Schur
complement makes \eqref{eq:445} equivalent to
\[
 \begin{pmatrix}\op P_{11}&\op P_{12}+\op X^*\\
                 \op P_{21}+\op X&\op P_{22}\end{pmatrix}
 -
 \begin{pmatrix}\op P_{13}\\\op P_{23}\end{pmatrix}
 \op P_{33}^{-1}
 \begin{pmatrix}\op P_{31}&\op P_{32}\end{pmatrix}\prec0.
\]
Due to \eqref{eq:446}, the negatives of the diagonal blocks are coercive.
The operator $\op X$ defined in \eqref{eq:447} renders the
off-diagonal block zero, so it is a PI solution.
\end{proof}

\begin{lemma}[Projection Lemma]
For bounded PI operators $\op A:\op H\to\op Y_A$,
$\op B:\op H\to\op Y_B$ and $\op P=\op P^*:\op H\to\op H$, consider
\begin{equation}\tag{4.4.8}\label{eq:448}
 \op A^*\op X\op B+\op B^*\op X^*\op A+\op P\prec0
\end{equation}
in the unstructured PI operator $\op X:\op Y_B\to\op Y_A$.
For the converse below, assume
\extra{that the block spaces used in the proof admit PI coordinates,
that the coordinate map $\op S$ is PI, and that
$\op S^*\op S$, $\op A_c^*\op A_c$ and $\op B_c^*\op B_c$
are coercive}.
The admissible PI class is closed under sums, products and adjoints;
zero-dimensional blocks are omitted.

Then \eqref{eq:448} has a solution if and only if
\begin{equation}\tag{4.4.9}\label{eq:449}
 \extra{\exists\,\delta>0:}\quad
 \op Ax=0\text{ or }\op Bx=0
 \ \Longrightarrow\
 \langle x,\op Px\rangle\le-\delta\|x\|^2.
\end{equation}
If $\op A_\perp$ and $\op B_\perp$ are isometric parameterizations
of $\ker\op A$ and $\ker\op B$, respectively,
\eqref{eq:449} is equivalent to
\begin{equation}\tag{4.4.10}\label{eq:4410}
 \op A_\perp^*\op P\op A_\perp\prec0
 \quad\text{and}\quad
 \op B_\perp^*\op P\op B_\perp\prec0.
\end{equation}
\end{lemma}

\begin{proof}
The proof of \emph{only if} is immediate. If $\op Ax=0$ or $\op Bx=0$,
the cross terms vanish, so \eqref{eq:448} implies \eqref{eq:449}.
This direction needs no additional coordinate or coercivity assumptions.

For proving \emph{if}, define
\[
 \begin{gathered}
 \op H_3=\ker\op A\cap\ker\op B,\qquad
 \op H_1=\ker\op A\ominus\op H_3,\\
 \op H_2=\ker\op B\ominus\op H_3,\qquad
 \op H_4=(\ker\op A+\ker\op B)^\perp.
 \end{gathered}
\]
Let $\op S:\bigoplus_{j=1}^4\op H_j\to\op H$ be
$\op S(x_1,x_2,x_3,x_4)=x_1+x_2+x_3+x_4$.
\extra{Coercivity of $\op S^*\op S$ makes $\op S$ bounded below
with closed range. Its range is dense by construction, so $\op S$
is boundedly invertible.}
Instead of \eqref{eq:448}, consider the equivalent inequality
\begin{equation}\tag{4.4.11}\label{eq:4411}
 (\op A\op S)^*\op X(\op B\op S)
 +(\op B\op S)^*\op X^*(\op A\op S)
 +\op S^*\op P\op S\prec0.
\end{equation}

Now $\op A\op S$ and $\op B\op S$ have the structure
$(0\ \op A_2\ 0\ \op A_4)$ and
$(\op B_1\ 0\ 0\ \op B_4)$.
The operators $\op A_c=(\op A_2\ \op A_4)$ and
$\op B_c=(\op B_1\ \op B_4)$ are injective.
\extra{Coercivity of $\op A_c^*\op A_c$ and $\op B_c^*\op B_c$}
gives the bounded PI left inverses
\[
 \op A_c^\dagger=(\op A_c^*\op A_c)^{-1}\op A_c^*,\qquad
 \op B_c^\dagger=(\op B_c^*\op B_c)^{-1}\op B_c^*.
\]
Thus the equation
\[
 (\op A\op S)^*\op X(\op B\op S)
 =
 \begin{pmatrix}0\\\op A_2^*\\0\\\op A_4^*\end{pmatrix}
 \op X\begin{pmatrix}\op B_1&0&0&\op B_4\end{pmatrix}
 =
 \begin{pmatrix}
 0&0&0&0\\
 \op Z_{21}&0&0&\op Z_{24}\\
 0&0&0&0\\
 \op Z_{41}&0&0&\op Z_{44}
 \end{pmatrix}
\]
has a PI solution for arbitrary bounded PI blocks
$\op Z_{21},\op Z_{24},\op Z_{41},\op Z_{44}$: take
\[
 \op X=(\op A_c^\dagger)^*
 \begin{pmatrix}\op Z_{21}&\op Z_{24}\\
                 \op Z_{41}&\op Z_{44}\end{pmatrix}
 \op B_c^\dagger.
\]
With $\op Q:=\op S^*\op P\op S$ partitioned accordingly,
\eqref{eq:4411} hence reads as
\begin{equation}\tag{4.4.12}\label{eq:4412}
\begin{pmatrix}
 \op Q_{11}&\op Q_{12}+\op Z_{21}^*&\op Q_{13}&\op Q_{14}+\op Z_{41}^*\\
 \op Q_{21}+\op Z_{21}&\op Q_{22}&\op Q_{23}&\op Q_{24}+\op Z_{24}\\
 \op Q_{31}&\op Q_{32}&\op Q_{33}&\op Q_{34}\\
 \op Q_{41}+\op Z_{41}&\op Q_{42}+\op Z_{24}^*&\op Q_{43}
   &\op Q_{44}+\op Z_{44}+\op Z_{44}^*
\end{pmatrix}\prec0.
\end{equation}
Since
\[
 \ker(\op A\op S)
 =\operatorname{im}\begin{pmatrix}I&0\\0&0\\0&I\\0&0\end{pmatrix},
 \qquad
 \ker(\op B\op S)
 =\operatorname{im}\begin{pmatrix}0&0\\I&0\\0&I\\0&0\end{pmatrix},
\]
the hypothesis \eqref{eq:449} amounts to
\[
 \begin{pmatrix}\op Q_{11}&\op Q_{13}\\
                 \op Q_{31}&\op Q_{33}\end{pmatrix}\prec0,
 \qquad
 \begin{pmatrix}\op Q_{22}&\op Q_{23}\\
                 \op Q_{32}&\op Q_{33}\end{pmatrix}\prec0.
\]
By Lemma~4.14, since $-\op Q_{33}$ is coercive and therefore
$\op Q_{33}^{-1}$ is PI, choose
\[
 \op Z_{21}=\op Q_{23}\op Q_{33}^{-1}\op Q_{31}-\op Q_{21}
\]
to make the negative of the leading $3\times3$ block in
\eqref{eq:4412} coercive. If $\op H_3=\{0\}$, take
$\op Z_{21}=-\op Q_{21}$ instead.

The blocks $\op Z_{41}$ and $\op Z_{24}$ can be arbitrary bounded
PI operators. After fixing them, write \eqref{eq:4412} as
\[
 \begin{pmatrix}
 \op R&\op F\\
 \op F^*&\op Q_{44}+\op Z_{44}+\op Z_{44}^*
 \end{pmatrix}\prec0,\qquad \op R\prec0.
\]
Choose $\op Z_{44}=-\alpha I$ with
\[
 2\alpha>\|\op Q_{44}\|+\|\op F\|^2\|\op R^{-1}\|.
\]
The negative of the Schur complement is coercive. This is the operator
counterpart of the original proof's application of Lemma~4.13.
Coercivity of $-\op R$ ensures that $\op R^{-1}$ is bounded and PI.
Pulling back through the boundedly invertible $\op S$
proves \eqref{eq:448}.
\end{proof}

\end{document}
